\section{Research design and scope}
\label{sec:research-design}

This paper is written as a method paper with an explicit empirical evaluation design. Its purpose is to define and justify a repeatable workflow for comparing situated and interview-based elicitation, tracing changes introduced by operator validation, and assessing whether validated fragments can be transformed into semantically coherent OSL representations. The method is intended for early-stage manufacturing projects where operator reasoning is important, but where Digital Twin requirements, models, and decision-support logic have not yet been formalised.

\subsection{Methodological position}

Notebook-to-OSL is a hybrid method rather than a new method family. It combines three layers. The elicitation layer is contextual and apprenticeship-based, with a conventional post-session interview used as a comparison source. The interpretation layer is cognitive-task-analysis-inspired: selected notes are examined for cues, difficult decisions, alternatives, expert--novice differences, exceptions, and escalation logic. The transformation layer is requirements- and MBSE-oriented: validated fragments are given traceability, status, rationale, and an explicit path toward formal representation.

Table~\ref{tab:method-family-comparison} summarises how the main literature streams contribute to the method.

\begin{table}[htbp]
\centering
\caption{Literature streams used to ground Notebook-to-OSL.}
\label{tab:method-family-comparison}
\begin{tabularx}{\textwidth}{p{0.25\textwidth}p{0.33\textwidth}X}
\toprule
Method family & Contribution to Notebook-to-OSL & Limitation if used alone \\
\midrule
Contextual inquiry and contextual design & Captures work in context, preserves operator phrasing, artefacts, flow, and interpretation checks \cite{beyer1995apprenticing,beyer1997contextual}. & Does not by itself establish what situated elicitation adds relative to a conventional interview. \\
Situated learning and cognitive apprenticeship & Justifies training as an elicitation setting where expert reasoning is demonstrated, coached, articulated, and reflected upon \cite{lave1991situated,collins1989cognitive}. & Provides theoretical framing but not a complete comparison or transformation protocol. \\
Cognitive task analysis & Surfaces cues, difficult judgments, expert--novice differences, anomalies, alternatives, and decision points \cite{crandall2006workingminds,klein1989critical,stanik2021acta}. & Can be intrusive during live work if applied as a full interview protocol. \\
Requirements engineering and SEBoK & Adds traceability, source, rationale, risk, status, ownership, and validation principles \cite{vanlamsweerde2009requirements,sebokStakeholder,sebokRequirements}. & Operates after stakeholder knowledge has been elicited and may conceal interpretation changes if provenance is weak. \\
MBSE and semantic DT modelling & Provides downstream targets such as SysML v2, controlled requirements, rules, and knowledge graphs \cite{mavin2009ears,omgSysml2024,karabulut2023ontologies,listl2024knowledge}. & Can over-formalise knowledge before it is reviewed and semantically qualified. \\
\bottomrule
\end{tabularx}
\end{table}

\subsection{Research-question alignment}

The three research questions define the evidence that the study must collect and the claims it may make. Table~\ref{tab:rq-alignment} maps each question to its data and analysis.

\begin{table}[htbp]
\centering
\caption{Alignment between research questions, evidence, and analysis.}
\label{tab:rq-alignment}
\begin{tabularx}{\textwidth}{p{0.09\textwidth}p{0.30\textwidth}p{0.28\textwidth}X}
\toprule
RQ & Required evidence & Analysis & Permitted conclusion \\
\midrule
RQ1 & Separate situated-training notes and post-session interview material covering the same task scope. & Compare which observations, contexts, interpretations, actions, exceptions, risks, and uncertainties occur in each source. & Whether and how the two elicitation settings produce different operator-strategy knowledge in the studied case. \\
RQ2 & Initial researcher annotations and operator-review records for the same fragments. & Trace accepted, corrected, narrowed, rejected, and sensitive interpretations and characterise the changes. & How operator review changes the researcher's interpretation and which categories are most vulnerable to misinterpretation. \\
RQ3 & Validated fragments, candidate OSL representations, transformation decisions, and unresolved gaps. & Assess traceability, semantic coherence, information loss, explicit incompleteness, and unsupported mappings. & Which validated fragments can be represented coherently, which remain incomplete, and where the representation requires refinement. \\
\bottomrule
\end{tabularx}
\end{table}

This alignment prevents the paper from treating the workflow, schema, and OSL example as independent contributions. They are research instruments used to answer the three questions.

\subsection{Design requirements}

The workflow should satisfy six requirements:

\begin{enumerate}[label=R\arabic*.]
    \item \textbf{Low friction during training.} The researcher must capture notes without interrupting the operator's explanation or the flow of the training situation.
    \item \textbf{Comparable sources.} Situated notes and post-session interview material must cover the same task scope and remain distinguishable during analysis.
    \item \textbf{Traceable interpretation.} Each structured fragment must remain linked to its raw source, the researcher's initial interpretation, operator-review outcome, and downstream representation.
    \item \textbf{Preservation of uncertainty.} The method must not force uncertain, context-dependent, alternative, or disputed observations into overly precise rules.
    \item \textbf{Validation before formalisation.} Operator or domain-expert review must occur before a fragment is treated as model-ready.
    \item \textbf{Assessable model-readiness.} The transformation must expose whether context, interpretation, response logic, uncertainty, provenance, and unresolved gaps can be represented coherently rather than merely producing syntactically plausible output.
\end{enumerate}

\subsection{Unit of analysis}

The primary unit of analysis is the \emph{operator-strategy fragment}. A fragment is derived from one or more raw notes or interview statements and relates an operational observation to context, interpretation, possible responses, rationale, exception, uncertainty, risk, responsibility, and validation state. The analysis also preserves the source-level distinction needed for RQ1, the version history needed for RQ2, and the transformation record needed for RQ3.

\subsection{Pilot study protocol}

A full empirical version of this paper should evaluate the method through a small but auditable pilot. The pilot should not aim to prove that a complete Digital Twin improves production performance. It should test the comparative, interpretive, and representational questions defined above.

The recommended pilot protocol is:

\begin{enumerate}
    \item Define one narrow task scope, such as one CNC monitoring task, recurring abnormality, or setup/inspection episode.
    \item Agree on consent and confidentiality boundaries, including what may be written, photographed, stored, anonymised, or excluded.
    \item Conduct a situated training session where the operator teaches the researcher in the real or simulated work setting.
    \item Capture low-friction digital field notes with identifiers, timestamps, raw text, setting, and optional task phase.
    \item Conduct a short post-session interview covering the same task scope without showing the situated notes as prompts.
    \item Keep situated and interview material separate, then code both with the same initial categories to answer RQ1.
    \item Create an initial researcher annotation for selected fragments and preserve that version.
    \item Apply CTA-style clarification where cues, alternatives, thresholds, or action logic remain ambiguous.
    \item Review the annotations with the operator and record each acceptance, correction, narrowing, rejection, or sensitivity decision to answer RQ2.
    \item Transform validated fragments into candidate OSL representations and record successful mappings, information loss, explicit gaps, and unsupported mappings to answer RQ3.
\end{enumerate}

\subsection{Empirical material and metrics}

Table~\ref{tab:empirical-material} summarises the minimum material needed for the pilot.

\begin{table}[htbp]
\centering
\caption{Minimum empirical material and its role in the research questions.}
\label{tab:empirical-material}
\begin{tabularx}{\textwidth}{p{0.27\textwidth}p{0.12\textwidth}X}
\toprule
Material & RQ & Purpose \\
\midrule
Situated training note log & RQ1 & Preserve knowledge captured while the operator demonstrates and explains work in context. \\
Post-session interview record & RQ1 & Provide a conventional elicitation comparison over the same task scope. \\
Initial annotation set & RQ2 & Preserve the researcher's interpretation before operator review. \\
Clarification and operator-review log & RQ2 & Record why interpretations were accepted, corrected, narrowed, rejected, or restricted. \\
Validated fragment set & RQ2--RQ3 & Establish the reviewed knowledge eligible for formal transformation. \\
Candidate OSL representations and mapping log & RQ3 & Show which semantic structures can be represented and which gaps remain. \\
Reflection log & All & Document researcher bias, missed context, methodological deviations, and confidentiality decisions. \\
\bottomrule
\end{tabularx}
\end{table}

For RQ1, the analysis should report source-exclusive and shared knowledge elements rather than only total note counts. For RQ2, it should report the number and type of interpretation changes and provide qualitative examples of consequential corrections. For RQ3, it should report the proportion of validated fragments that can be represented with traceability and semantic coherence, the proportion that remain explicitly incomplete, and the recurring reasons why transformation fails or loses information. These measures are descriptive and case-bounded; they do not establish general superiority, operational correctness, or safety.

\subsection{Scope and boundaries}

The paper does not claim to deliver a complete Operator Strategy Language or a complete Digital Twin. It focuses on the intermediate research problem between situated operator learning and formal modelling. RQ1 is limited to the elicitation settings studied. RQ2 evaluates interpretation change, not universal operator agreement. RQ3 evaluates representation adequacy for the validated fragments, not runtime matching, automatic recommendation, operational correctness, or safety.

The tool is treated as a research instrument that supports capture, versioning, review, and transformation. OSL is treated as the formal target needed to examine RQ3. Neither the tool nor the language is validated merely because the workflow can produce an artefact.
